%% S7: no closed S-term reduces to a term that contains it (Waldmann's Cor. 17), constructively.
%% DRAFT 2026-11-23 (skinock).  Neutral amsart shell, as S1--S6.
%% Paper-to-Lean map (lean-ski, modules Ski/HeadRec.lean since v1.8.0 and Ski/HeadRecF.lean since v1.15.0; all in #audit_clean):
%%   lem:label ...... HeadRec.acc_all (acc_app, acc_app_aux; relation Rel, labelled step LStep)
%%   lem:bounded .... HeadRec.bsn (tag, Cons, lift, rd_bstep, cons_tag)
%%   thm:main ....... HeadRec.no_ground_loop (sub_replay)
%%   cor:nhr ........ HeadRec.no_head_recurrence (sub_apps)
%%   thm:f .......... HeadRec.theorem_H (Ski/HeadRecF.lean, v1.15.0; FStep, confluent, nf_unique, fstep_tag, sub_lift, le_sims, h_core)
\documentclass[11pt]{amsart}
\usepackage[T1]{fontenc}
\usepackage{lmodern}
\usepackage{amsmath,amssymb,amsthm}
\usepackage[numbers,sort]{natbib}
\usepackage[hidelinks]{hyperref}
\usepackage{bm}

\newtheorem{theorem}{Theorem}
\newtheorem{lemma}[theorem]{Lemma}
\newtheorem{corollary}[theorem]{Corollary}
\newtheorem{conjecture}[theorem]{Conjecture}

\newcommand{\Ter}{\mathcal S}
\newcommand{\red}{\twoheadrightarrow}
\newcommand{\conv}{=_\beta}
\newcommand{\rd}{\operatorname{rd}}
\newcommand{\lab}[1]{\mathbin{\cdot_{#1}}}

\title{No Ground Loops for the Combinator $\bm{S}$, Constructively}
\author{Claude Opus \& Fable with N.\ E.\ Davis}
\address{ZeroAttest, Champaign, Illinois, USA}
\email{NDavis@ZeroAttest.com}
\urladdr{https://orcid.org/0000-0002-1131-8319}
\date{\today}

\begin{document}
\begin{abstract}
Waldmann proved that no closed term built from the combinator $S$ reduces, in one or more steps,
to a term that contains it. We give a constructive proof and mechanize it in Lean~4. We
prove termination of a variant of Waldmann's labelled reduction by direct induction on accessibility instead
of a recursive path order. The theorem follows by replaying a loop inside its own result. As a consequence, no closed $S$-term head-reduces to itself applied to arguments. More generally, if a term reaches $u$ without contracting an argument, for instance by head steps, then it reduces to no term $C[u]$ with $C\ne[\,]$ a context.
\end{abstract}

\keywords{combinatory logic, combinator $S$, termination, accessibility, Lean}
\maketitle

\section{Introduction}

$\Ter$ is the set of closed terms built from $S$ by application, with $S\,x\,y\,z\to x\,z\,(y\,z)$ at
any position; $\to^+$ is the transitive closure and $\red$ the reflexive--transitive closure. A
\emph{head step} contracts the redex at the head of the left spine. A \emph{ground loop} is a reduction $t\to^+u$ in which $t$ occurs in $u$. Waldmann showed that
there are none~\citep[Cor.~17]{waldmann-ic}. His proof labels the terms and shows that the labelled
system terminates by a recursive path order~\citep[Thm.~8]{waldmann-ic}. We use a variant of his labelling,
and replace the order and the infinite reduction of~\citep[Cor.~17]{waldmann-ic} by a direct
induction on accessibility, which avoids both; the whole argument is checked in Lean.

\begin{theorem}\label{thm:main}
No $t\in\Ter$ has $t\to^+u$ with $t$ a subterm of $u$ (possibly $u=t$).
\end{theorem}

\begin{corollary}\label{cor:nhr}
No $t\in\Ter$ head-reduces in one or more steps to $t\,X_1\cdots X_k$, for any $k\ge0$.
\end{corollary}

\begin{proof}
Head steps are steps, and $t$ occurs in $t\,X_1\cdots X_k$.
\end{proof}

\section{Labelled reduction}

The \emph{right depth} is $\rd(S)=0$ and $\rd(f\,x)=\rd(x)+1$. A \emph{labelled term} is $S$ or
$a\lab k b$ with labelled terms $a,b$ and $k\in\mathbb N$. A \emph{labelling} of $t\in\Ter$ is a labelled term that becomes $t$ when its labels are erased. For $N\in\mathbb N$, an \emph{$N$-step}
rewrites, at any position,
\[
  ((S\lab i x)\lab j y)\lab k z\;\to\;(x\lab k z)\lab{k+1}(y\lab k z)\qquad\text{if }k<N.
\]
Write $a\prec b$ if $b$ makes an $N$-step to $a$, if $a$ is the function part or the argument of
$b$, or if $b=(l\lab i a)\lab k r$. The \emph{accessible} terms are defined inductively: a term is accessible if all its
$\prec$-predecessors are.

\begin{lemma}[{cf.\ \citep[Thm.~8]{waldmann-ic}}]\label{lem:label}
For every $N$, every labelled term is accessible; in particular $N$-steps terminate.
\end{lemma}

\begin{proof}
We show that $u\lab k r$ is accessible whenever $u$ and $r$ are, by induction on $\max(N-k,0)$, and
within it by induction on the accessibility of $u$, then of $r$. The predecessors of $u\lab k r$ are
$u$, $r$, the argument $x$ of $u=l\lab i x$ (which precedes $u$), and the results of $N$-steps. A
step inside $u$ or inside $r$ gives a term covered by the inner hypotheses. A step at the root
needs $k<N$ and $u=(S\lab i x)\lab j y$, and gives $(x\lab k r)\lab{k+1}(y\lab k r)$. Here $x$ and $y$
precede $u$, so the hypothesis for $u$ makes $x\lab k r$ and $y\lab k r$ accessible, and the outer
hypothesis for $k+1$, where $N-(k+1)<N-k$, makes the result accessible. Since $S$ has no
predecessors, induction on the structure of labelled terms finishes the proof.
\end{proof}

\begin{lemma}[{cf.\ \citep[Thm.~16]{waldmann-ic}}]\label{lem:bounded}
For every $N$, every $t\in\Ter$ is accessible for reduction restricted to redexes $S\,f\,g\,x$
with $\rd(x)<N$; in particular this reduction terminates.
\end{lemma}

\begin{proof}
Call a labelling of $t$ \emph{consistent} if every label is at most the right depth of the argument
below it; labelling each node $f\,x$ with $\rd(x)$ is consistent. A restricted step never lowers the
right depth of a term: at the root
$\rd(S\,f\,g\,x)=\rd(x)+1<\rd(x)+2=\rd(f\,x\,(g\,x))$, and at other positions by induction on
the term. So the labels away from the contracted redex stay consistent. At the redex, the label
$k\le\rd(x)<N$ lets the $N$-step fire, and
the new labels satisfy $k\le\rd(x)$ and $k+1\le\rd(g\,x)$. So for every consistent labelling $T$ of $t$ and restricted step $t\to t'$ there is an
$N$-step $T\to T'$ with $T'$ a consistent labelling of $t'$. By induction on the accessibility of
$T$ (Lemma~\ref{lem:label}), $t$ is accessible; take $T$ the right-depth labelling.
\end{proof}

\section{Proof of Theorem~\ref{thm:main}}

Suppose $t\to t'\red u$ with $t$ a subterm of $u$, and let $N$ exceed the right depths of the
arguments of the finitely many redexes contracted. Then the reduction is restricted in the sense of
Lemma~\ref{lem:bounded}. If $t$ occurs in a term $w$, the same steps at that occurrence (the
\emph{replay}) give
restricted reductions $w\to w'\red w''$ with $t$ occurring in $w''$.

We show, by induction on the accessibility of $w$ for restricted steps, that no restricted
reduction $w\red v$ has $t$ occurring in $v$. If $w\to w_0\red v$, apply the hypothesis to $w_0$. If
$v=w$, then $t$ occurs in $w$, and the replay gives $w\to w'\red w''$ with $t$ in $w''$, contradicting the
hypothesis for $w'$. Since $t$ is accessible by Lemma~\ref{lem:bounded}, taking $w=v=t$ gives a
contradiction.\qed

\section{Reducts reached without contracting an argument}

A \emph{context} $C$ is a term with one hole $[\,]$, and $C[u]$ fills the hole with $u$. A step is
an \emph{f-step} if its redex is the whole term or the function part of its parent, that is, not an
argument (it may lie inside one); head steps are f-steps. Write $t\red_f u$ for a reduction by f-steps.

\begin{theorem}\label{thm:f}
If $t\in\Ter$, $t\red_f u$ and $t\red C[u]$ for a context $C$, then $C=[\,]$.
\end{theorem}

\noindent Taking $u=t$, with no f-steps, gives the case of Theorem~\ref{thm:main} in which $t$ is a
proper subterm of the reduct.

\begin{proof}
Let $t_{\rd}$ label each node $f\,x$ of $t$ with $\rd(x)$, write $T\le T'$ if two labellings of one
term are pointwise ordered, and let $|T|$ count the occurrences of $S$. Fix $N$ above the right
depths of the arguments of the redexes contracted in both reductions. $N$-steps are orthogonal: their left-hand sides are linear, the
proper subpatterns $(S\lab i x)\lab j y$ and $S\lab i x$ of a redex are not redexes, and the
condition $k<N$ reads only the label at the root of the redex. Orthogonal systems are
confluent~\citep[Ch.~4]{terese}, so by Lemma~\ref{lem:label} every labelled term $T$ has a unique
$N$-normal form $\hat T$ (a reduct with no $N$-step); and $|\hat T|\ge|T|$, since an $N$-step deletes one $S$ and copies $z$,
which contains one.

An f-step $t\to t'$ gives an $N$-step $t_{\rd}\to t'_{\rd}$: the new labels are $\rd(x)$, $\rd(x)$ and
$\rd(x)+1=\rd(g\,x)$ for the redex $S\,f\,g\,x$, and the right depth of an argument depends only on
its right spine, whose nodes are arguments, and an f-step contracts none of them. So $t_{\rd}$ reduces
by $N$-steps to $u_{\rd}$ and, by the proof of Lemma~\ref{lem:bounded}, to a consistent labelling
$C'[R]$ of $C[u]$. Consistency compares a label only with the argument below it, so $R$ is a
consistent labelling of $u$, that is, $R\le u_{\rd}$.

If $T\le T'$ then $|\hat T|\ge|\hat T'|$: an $N$-step of $T'$ fires at the same redex of $T$ and
keeps the order, so $T$ reduces to some $V\le\hat T'$, and $|\hat T|=|\hat V|\ge|V|=|\hat T'|$.
Now $\hat u_{\rd}=\widehat{C'[R]}=\widehat{C'[\hat R]}$. Suppose $C\ne[\,]$. Then $C$ contains an
$S$ outside the hole, and $|\hat u_{\rd}|=|\widehat{C'[\hat R]}|\ge|C'[\hat R]|>|\hat R|\ge|\hat
u_{\rd}|$, which is absurd.
\end{proof}

\section{Closing}

Let $\conv$ be convertibility, the equivalence closure of $\to$. Theorem~\ref{thm:main} excludes $t\to^+E[t]$. Excluding $t\conv E[t]$ for $E\ne[\,]$ is stronger: in
the $\lambda$-calculus, with Curry's fixed-point combinator $Y$, the term $Y\,f$ is convertible with
$f\,(Y\,f)$, although $Y\,f$ does not reduce in one or more steps to a term containing $Y\,f$.

\begin{conjecture}
No $t\in\Ter$ and context $E\ne[\,]$ satisfy $t\conv E[t]$.
\end{conjecture}

\noindent By Church--Rosser, the conjecture says that $t$ and $E[t]$ have no common reduct. It holds
when $t$ has a normal form $n$. Then $E[n]\conv n$, so $E[n]\red n$. But a step
$S\,x\,y\,z\to x\,z\,(y\,z)$ never lowers the number of occurrences of $S$, since $z$ contains one,
and $E[n]$ has more occurrences than $n$. The conjecture would strengthen Statman's theorem that $S$ alone
defines no fixed-point combinator~\citep[Rem.~15]{dots}: such a $Y$ would give $Y\,S\conv S\,(Y\,S)$.

\section*{Data availability}
Lemmas~\ref{lem:label} and~\ref{lem:bounded}, Theorems~\ref{thm:main} and~\ref{thm:f}, and
Corollary~\ref{cor:nhr} are proved in the Lean~4~\citep{lean4} development
\texttt{lean-ski}~\citep{leanski}, modules \texttt{HeadRec} and \texttt{HeadRecF}, using only the axioms
\texttt{propext} and \texttt{Quot.sound}.

\bibliographystyle{plainnat}
\bibliography{refs}
\end{document}
